% Template for web background cards
% Parameters:
%   \bgwidth - Card width
%   \bgheight - Card height
%   \bgcolourbase - Base colour for background (e.g._light,_dark)
%   \colourone - First accent colour
%   \colourtwo - Second accent colour
%   \opacityset - Opacity multiplier (e.g., 0.5 for light, 0.8 for dark)
%   \browserx - X position for browser window
%   \browsery - Y position for browser window
%   \iconsx - X position for responsive icons
%   \iconsy - Y position for responsive icons
%   \codex - X position for code snippet
%   \codey - Y position for code snippet
%   \palettex - X position for colour palette
%   \palettey - Y position for colour palette
%   \gridx - X position for grid
%   \gridy - Y position for grid
%   \labelonex - X position for first label
%   \labeloney - Y position for first label
%   \labeltwox - X position for second label
%   \labeltwoy - Y position for second label
%   \labelthreex - X position for third label
%   \labelthreey - Y position for third label
%   \lineoneax - X start for first decorative line
%   \lineoneay - Y position for first decorative line
%   \lineonebx - X end for first decorative line
%   \linetwoax - X start for second decorative line
%   \linetwoay - Y position for second decorative line
%   \linetwobx - X end for second decorative line
%   \cornerx - X position for corner accent
%   \cornery - Y position for corner accent
%   \corneryb - Y bottom for corner accent
%   \cornerxb - X right for corner accent

\begin{tikzpicture}
  % Define the card background
  \fill[\bgcolourbase] (0,0) rectangle (\bgwidth,\bgheight);
  
  % Browser window mockup
  \begin{scope}[shift={(\browserx,\browsery)}]
    % Browser chrome
    \draw[\colourone, opacity=\opacityone, rounded corners=2pt] (0,0) rectangle (3.0,1.5);
    \fill[\colourone, opacity=\opacitychrome, rounded corners=2pt] (0,1.3) rectangle (3.0,1.5);
    % Browser dots
    \fill[\colourtwo, opacity=\opacityone] (0.15,1.37) circle (0.04);
    \fill[\colourtwo, opacity=\opacityone] (0.30,1.37) circle (0.04);
    \fill[\colourtwo, opacity=\opacityone] (0.45,1.37) circle (0.04);
    % Content area with simple layout
    \fill[\colourtwo, opacity=\opacitycontent] (0.2,0.9) rectangle (2.8,1.2);
    \fill[\colourone, opacity=\opacitycontent] (0.2,0.2) rectangle (1.3,0.8);
    \fill[\colourtwo, opacity=\opacitycontent] (1.5,0.2) rectangle (2.8,0.8);
  \end{scope}
  
  % Responsive design icons (mobile, tablet, desktop)
  \begin{scope}[shift={(\iconsx,\iconsy)}, scale=0.7]
    % Mobile
    \draw[\colourtwo, thick, opacity=\opacityicons, rounded corners=1pt] (0,0) rectangle (0.4,0.8);
    \fill[\colourtwo, opacity=\opacityiconfill] (0.05,0.15) rectangle (0.35,0.65);
    % Tablet
    \draw[\colourone, thick, opacity=\opacityicons, rounded corners=1pt] (0.7,0.1) rectangle (1.3,0.7);
    \fill[\colourone, opacity=\opacityiconfill] (0.8,0.2) rectangle (1.2,0.6);
    % Desktop
    \draw[\colourtwo, thick, opacity=\opacityicons] (1.7,0.2) rectangle (2.6,0.7);
    \draw[\colourtwo, thick, opacity=\opacityicons] (2.0,0.2) -- (2.3,0.05) -- (2.3,0.05);
    \fill[\colourtwo, opacity=\opacityiconfill] (1.75,0.25) rectangle (2.55,0.65);
  \end{scope}
  
  % Code snippet element
  \begin{scope}[shift={(\codex,\codey)}]
    \draw[\colourone, opacity=\opacitycodebg, rounded corners=2pt] (0,0) rectangle (3.5,1.2);
    \node[\colourone, font=\tiny\ttfamily, opacity=\opacitycode, align=left, anchor=west] at (0.15,0.85) 
      {\textless div class="content"\textgreater};
    \node[\colourtwo, font=\tiny\ttfamily, opacity=\opacitycode, align=left, anchor=west] at (0.3,0.55) 
      {\textless h1\textgreater Hello World\textless/h1\textgreater};
    \node[\colourone, font=\tiny\ttfamily, opacity=\opacitycode, align=left, anchor=west] at (0.15,0.25) 
      {\textless/div\textgreater};
  \end{scope}
  
  % Colour palette squares
  \begin{scope}[shift={(\palettex,\palettey)}]
    \fill[\colourtwo, opacity=\opacitypaletteone] (0,0) rectangle (0.4,0.4);
    \fill[\colourone, opacity=\opacitypaletteone] (0.5,0) rectangle (0.9,0.4);
    \fill[\colourtwo, opacity=\opacitypalettetwo] (1.0,0) rectangle (1.4,0.4);
    \fill[\colourone, opacity=\opacitypalettetwo] (1.5,0) rectangle (1.9,0.4);
  \end{scope}
  
  % Grid layout representation
  \begin{scope}[shift={(\gridx,\gridy)}, opacity=\opacitygrid]
    \draw[\gridcolour, thick] (0,0) grid[step=0.4] (1.6,1.2);
  \end{scope}
  
  % Friendly labels
  \node[\colourone, font=\small\sffamily, opacity=\opacitylabels] at (\labelonex, \labeloney) {Design};
  \node[\colourtwo, font=\small\sffamily, opacity=\opacitylabels] at (\labeltwox, \labeltwoy) {Responsive};
  \node[\colourone, font=\small\sffamily, opacity=\opacitylabels] at (\labelthreex, \labelthreey) {Creative};
  
  % Subtle decorative elements
  \draw[\colourtwo, thick, opacity=\opacitylines] (\lineoneax,\lineoneay) -- (\lineonebx,\lineoneay);
  \draw[\colourone, thick, opacity=\opacitylines] (\linetwoax,\linetwoay) -- (\linetwobx,\linetwoay);
  
  % Small corner accent
  \draw[\colourtwo, line width=2pt, opacity=\opacitycorner] (\cornerx,\cornery) -- (\cornerx,\corneryb);
  \draw[\colourtwo, line width=2pt, opacity=\opacitycorner] (\cornerx,\cornery) -- (\cornerxb,\cornery);
  
\end{tikzpicture}
